\section{Step 5 Read the two bytes of input A}\label{pcm-s05}

Load offsets zero and one from \texttt{a0} with \texttt{lbu}. Each load yields an unsigned byte. Two separate byte operations permit a sample to begin at an odd address and make the logical two-byte boundary explicit. Do not interpret the high byte as a signed sample by itself: its sign bit belongs to the complete halfword.

\Needspace{12\baselineskip}
\subsection{Write and inspect}
Replace the temporary branch and nop in \texttt{load\_a}. Keep the surrounding labels and region markers.

\Needspace{5\baselineskip}
\MARSCodeLabel{MIPS assembly}
\begin{minted}[samepage=true]{mars-idea}
    lbu $t0, 0($a0)
    lbu $t1, 1($a0)
\end{minted}
Use recipe \texttt{s05-1} from the companion construction inputs. Pause at \texttt{mix\_ready} before comparing this partial state.

\subsection{Expected observation}
For A=-1000, the bytes are hexadecimal 18 fc: \texttt{t0}=24 and \texttt{t1}=252.

\Needspace{11\baselineskip}
\subsection{Check your understanding}
Why do the two bytes not each represent a separate sample?

\begin{itemize}
  \item \textbf{A} Every byte is a complete sixteen-bit value.
  \item \textbf{B} Together they encode one signed halfword.
  \item \textbf{C} The loads advance the pointer automatically.
\end{itemize}
